% AI-623 Deep Vision Language Models — Final Report
% Group 13: Adnan, Hatim, Adnan
%
% Build:
%   pdflatex group13_v2.tex && bibtex group13_v2 && pdflatex group13_v2 && pdflatex group13_v2
%
% Requires: neurips_2024.sty in the same directory.
% Figures expected under ../outputs/figures/ as PDFs.

\documentclass{article}

% NeurIPS 2024 style.
\usepackage[final]{neurips_2024}

\usepackage[utf8]{inputenc}
\usepackage[T1]{fontenc}
\usepackage{url}
\usepackage{booktabs}
\usepackage{amsfonts}
\usepackage{amsmath}
\usepackage{nicefrac}
\usepackage{microtype}
\usepackage{xcolor}
\usepackage{graphicx}
\usepackage{subcaption}
\usepackage{hyperref}
\usepackage{makecell}
\usepackage{multirow}
\usepackage{enumitem}

\graphicspath{{../outputs/figures/}}

\title{Pixel-Space Posterior Sampling for Image Deblurring:\\
       A Controlled Comparison of DPS and FlowDPS-Style Guidance\\
       on CelebA-HQ-256}

\author{%
  Eeman Adnan \\
  \texttt{24280022@lums.edu.pk} \\
  \And
  Muhammad Aslam Adnan \\
  \texttt{25280067@lums.edu.pk} \\
  \And
  Mufaddal Hatim \\
  \texttt{25280084@lums.edu.pk} \\
}

\begin{document}

\maketitle

\begin{abstract}
Image reconstruction from degraded observations is a fundamental inverse
problem in computational imaging. \emph{Diffusion Posterior Sampling}
(DPS)~\cite{chung2023dps} leverages a learned diffusion prior plus a
measurement-likelihood gradient to solve such problems with high
reconstruction fidelity, at the cost of hundreds of neural function
evaluations (NFE). \emph{Flow Matching} models~\cite{lipman2023flow,liu2023rectified}
offer a deterministic, ODE-based alternative that has been argued to
enable comparable quality with fewer steps. In this work we conduct a
controlled comparison of \textbf{pixel-space DPS} on a CelebA-HQ-256 DDPM
against a \textbf{pixel-space re-implementation of FlowDPS}~\cite{kim2025flowdps},
whose official code is latent-only (Stable Diffusion 3), on top of the
Liu~et~al.\ 2023 Rectified Flow CelebA-HQ-256 checkpoint. We evaluate
both methods on Gaussian deblurring across $3 \times 2 \times 3 = 18$
conditions $\times\,50$ images $= 1800$ reconstructions, and study
operator-mismatch robustness using motion-blur measurements while the
sampler's likelihood assumes Gaussian (1200 additional reconstructions).
We further include Wiener-filter and DnCNN$+$PnP-ADMM baselines (600
reconstructions). On the matched-Gaussian grid Pixel-DPS dominates the
comparison, with a PSNR advantage of $+1.4$ to $+4.2$\,dB over our
FlowDPS reimplementation, and FlowDPS saturates with NFE while Pixel-DPS
keeps improving. Under operator mismatch, however, FlowDPS-on-RF is
\emph{relatively more robust} (degrades by $\sim$2.5\,dB versus
$\sim$5\,dB for Pixel-DPS), reversing the matched-condition comparison.
All four methods comfortably beat Wiener; DnCNN$+$PnP-ADMM beats our
FlowDPS reimplementation in $5/6$ noisy cells while running $\sim$60$\times$
faster than DPS. We additionally evaluate \emph{eight} publication-improvement
candidates spanning EMA-aligned weight loading, time-varying $\zeta$
schedules, spectral-aware likelihood reweighting, $\Pi$-GDM-style
covariance correction, and a 2nd-order Heun integrator. Only one
(\texttt{flowdps\_rf\_v2}: EMA load $+$ ramp-$\zeta$) clears a strict
Bonferroni-corrected significance gate, adding $+0.77$\,dB to FlowDPS-on-RF;
the remaining seven are documented as failed ablations with diagnoses.
The findings give concrete guidance for choosing between these methods
under known versus uncertain forward operators.
\end{abstract}

% =====================================================================
\section{Introduction}\label{sec:intro}

Image reconstruction from degraded observations is a fundamental problem
in computer vision and computational imaging. It is typically formulated
as an inverse problem,
\begin{equation}
  y = A x + \epsilon, \qquad \epsilon \sim \mathcal{N}(0, \sigma_n^2 I),
\end{equation}
where $A$ is a known or partially known degradation operator (e.g.\
Gaussian blur), and the goal is to recover the unknown clean image $x$
given only $y$ and a prior over natural images. The problem is inherently
ill-posed; multiple plausible $x$ can produce the same observation, so a
strong prior is essential. Modern approaches replace classical
handcrafted priors with learned generative priors such as Denoising
Diffusion Probabilistic Models (DDPMs)~\cite{ho2020ddpm} and Flow Matching
models~\cite{lipman2023flow,liu2023rectified}. These priors combined with
explicit measurement guidance form the basis of methods like
DPS~\cite{chung2023dps}, DDRM~\cite{kawar2022ddrm}, and the more recent
FlowDPS~\cite{kim2025flowdps}.

This paper conducts a controlled empirical comparison of two posterior
sampling paradigms for image deblurring on faces:
\begin{enumerate}[leftmargin=*,itemsep=2pt]
  \item \textbf{Pixel-DPS}: the original DPS sampler from
    Chung~et~al.~\cite{chung2023dps}, on a CelebA-HQ-256 DDPM checkpoint
    (\texttt{google/ddpm-ema-celebahq-256}).
  \item \textbf{FlowDPS-on-RF}: a from-scratch pixel-space
    reimplementation of the Kim~et~al.~2025 FlowDPS algorithm on top of
    the Liu~et~al.~2023 CelebA-HQ-256 Rectified Flow
    checkpoint~\cite{liu2023rectified}. The official Kim 2025 codebase is
    Stable~Diffusion~3-only; our reimplementation exists so we can do a
    true pixel-space comparison against Pixel-DPS at the same resolution
    and on the same dataset.
\end{enumerate}

Both methods therefore share the \emph{same dataset, same resolution, and
same NCSN++/U-Net architectural family}, so any quality or efficiency
difference can be cleanly attributed to the choice of posterior-sampling
paradigm rather than to incidental differences in the underlying prior.
We additionally evaluate two classical baselines (Wiener,
DnCNN$+$PnP-ADMM) and one operator-mismatch robustness study (true
measurement is motion blur; sampler's likelihood assumes Gaussian).

The study is guided by four research questions:
\begin{enumerate}[leftmargin=*,itemsep=2pt]
  \item \textbf{Quality.} Can a flow-based posterior sampler (FlowDPS on
    Rectified Flow) match diffusion-based DPS in reconstruction
    fidelity (PSNR / SSIM / LPIPS) at matched NFE budgets?
  \item \textbf{Efficiency.} How does per-image wall-clock scale with
    NFE for each method, and is there a regime where the deterministic
    flow trajectory yields a measurable speed advantage?
  \item \textbf{Robustness.} Under forward-model misspecification (true
    measurement is motion blur, sampler's likelihood assumes Gaussian),
    which method degrades more gracefully?
  \item \textbf{Baselines.} How do the diffusion / flow posterior
    samplers compare to classical (Wiener) and supervised
    (DnCNN$+$PnP-ADMM) reconstruction baselines on identical inputs?
\end{enumerate}

The contributions of this work are:
\begin{itemize}[leftmargin=*,itemsep=2pt]
  \item A pixel-space, fair-comparison evaluation of two posterior
    samplers (DPS and FlowDPS-style guidance) on a common dataset and
    architectural family.
  \item Empirical evidence that, on this setup, Pixel-DPS dominates
    FlowDPS-on-RF in the matched-Gaussian setting by 1.4--4.2\,dB PSNR,
    while FlowDPS-on-RF degrades less under operator mismatch ---
    reversing the matched-condition ordering.
  \item A four-method (Wiener, DnCNN$+$PnP, FlowDPS-on-RF, Pixel-DPS)
    Pareto-frontier characterization across three orders of magnitude
    of per-image compute.
  \item A rigorous \emph{ablation study} of eight candidate
    improvements to either sampler (EMA-aligned weight loading,
    time-varying $\zeta$ schedules, spectral-domain residual
    reweighting, $\Pi$-GDM-style covariance correction, and a Heun
    2nd-order integrator), evaluated against a strict
    Bonferroni-corrected significance gate. We report a single
    confirmed win ($+0.77$\,dB) together with seven diagnosed
    failures; the negative results --- including a subtle DPS
    L2-norm vs.\ $\Pi$-GDM L2-squared convention mismatch and a
    guidance-frequency vs.\ integration-accuracy trade-off in Heun ---
    are instructive in their own right.
  \item A reproducible open-source pipeline at
    \url{https://github.com/adnan821/dps_flowdps}, including the 3{,}600
    measured rows, all per-experiment logs, and the full LaTeX source
    for this report.
\end{itemize}

% =====================================================================
\section{Related Work}\label{sec:related}


Classical inverse-problem solvers rely on handcrafted regularizers and
analytical priors. The Wiener filter solves the linear MMSE problem in
the Fourier domain for Gaussian degradations and remains a fast,
interpretable baseline. Deep-learning denoiser priors enabled
Plug-and-Play~(PnP) schemes such as PnP-ADMM~\cite{chan2017pnp,zhang2017dncnn}
that swap a proximal-step argmin for a forward pass through a pretrained
denoiser network (e.g.~DnCNN, DRUNet). Diffusion-based posterior sampling
generalized this idea by injecting a measurement-consistency term into
the reverse diffusion process: DDRM~\cite{kawar2022ddrm} uses spectral
projection in the SVD basis of $A$, and DPS~\cite{chung2023dps} uses a
likelihood-gradient term computed via Tweedie's clean-image estimate.

Flow Matching and Rectified Flow~\cite{lipman2023flow,liu2023rectified}
recently emerged as deterministic alternatives to DDPMs, learning a
velocity field whose ODE integration straightens the noise-to-data path
and can in principle achieve comparable image quality with fewer
function evaluations. FlowDPS~\cite{kim2025flowdps} adapts DPS-style
likelihood guidance to flow models by injecting a likelihood-gradient
correction into the deterministic velocity field. The official FlowDPS
implementation is built on the SD3 latent transformer flow model and
operates at $768 \times 768$ in latent space.

Beyond DPS-style likelihood-gradient guidance, several recent works
explore alternative posterior-sampling regimes for diffusion- and
flow-based priors. Meanti et~al.~\cite{meanti2025dim} cast the
inverse problem as distribution matching between a parameterized
sampler and a target conditioned on the measurement, avoiding explicit
likelihood gradients. Dou et~al.~\cite{dou2026particle} propose a
\emph{gradient-free} constrained-particle approach that solves
diffusion inverse problems using only forward passes through the
score / flow network, side-stepping the cost of backpropagating through
the model that DPS-style methods incur. Dasgupta et~al.~\cite{dasgupta2026pcfm}
study \emph{physics-constrained} conditional flow matching for inverse
problems with hard physical constraints (e.g.\ PDE residuals), where
posterior-sample physical validity matters as much as image quality.
Our work uses gradient-based DPS / FlowDPS guidance; we leave a
gradient-free or physics-constrained variant to future work.

To enable a like-for-like pixel-space comparison against Pixel-DPS, we
reimplement the FlowDPS likelihood-corrected velocity update,
\begin{equation}
  v_\text{eff}(x_t, t) = v_\theta(x_t, t) - \zeta \nabla_{x_t}
    \| y - A( \hat z_1(x_t, t) ) \|_2,
\end{equation}
on top of the Liu~2023 Rectified Flow CelebA-HQ-256 NCSN++ velocity
model, where the Tweedie clean-image identity for rectified flow is
$\hat z_1(x_t, t) = x_t + (1 - t)\,v_\theta(x_t, t)$.

% =====================================================================
\section{Methodology}\label{sec:method}

\paragraph{Forward model.}
$y = A x + \epsilon$, with $A$ either a Gaussian blur of standard
deviation $\sigma_b \in \{1.5, 3.0, 5.0\}$ (main grid) or a motion blur
$\text{motion}(L,\theta)$ with $L \in \{15, 25, 35\}$ pixels and
$\theta \in \{0^\circ, 45^\circ\}$ (robustness study), and
$\epsilon \sim \mathcal{N}(0, \sigma_n^2 I)$ with
$\sigma_n \in \{0.0, 0.05\}$.

\paragraph{Pixel-DPS.}
We adopt the DDPM reverse process with DDIM sampling. At each reverse
step $t$, the model predicts $\epsilon_\theta(x_t, t)$, from which we
derive the Tweedie clean-image estimate
\begin{equation}
  \hat x_0(x_t) = \tfrac{x_t - \sqrt{1 - \bar\alpha_t}\,\epsilon_\theta(x_t, t)}{\sqrt{\bar\alpha_t}}.
\end{equation}
The DDIM-style reverse step is then corrected with a likelihood gradient:
\begin{equation}
  x_{t-1} = x_{t-1}^{\text{DDIM}} - \zeta \nabla_{x_t} \| y - A(\hat x_0(x_t)) \|_2.
\end{equation}
Two implementation choices matter substantially. First, we use the
\textbf{L2 norm} (not the squared L2) of the residual; squaring the
residual and dividing by $\sigma_n^2$ blows up the gradient and prevents
convergence (verified empirically in our EXP-001 ablation). Second, we do
\textbf{not clamp} $\hat x_0$ inside the gradient path; clamping zeroes
the gradient for any pixel outside $[-1, 1]$, which kills DPS at early
timesteps when $\hat x_0$ is necessarily far from the data manifold.
Clamping is applied only at the final returned tensor.

\paragraph{FlowDPS-on-RF.}
We use the Liu~2023 CelebA-HQ-256 1-Rectified Flow checkpoint, which
provides a learned velocity field $v_\theta(x_t, t)$ over continuous time
$t \in [0,1]$, with $x_0 \sim \mathcal{N}(0, I)$ at $t = 0$ and $x \sim p_\text{data}$
at $t = 1$. Sampling proceeds by forward Euler integration over a uniform
$N$-step grid in $[\epsilon, 1 - \epsilon]$, with FlowDPS guidance:
\begin{equation}
  v_\text{eff}(x_t, t) = v_\theta(x_t, t) - \zeta \nabla_{x_t} \| y - A(\hat z_1(x_t, t)) \|_2,
  \qquad
  x_{t + dt} = x_t + dt \cdot v_\text{eff}(x_t, t).
\end{equation}
The clean-image estimate $\hat z_1 = x_t + (1-t) v_\theta(x_t, t)$ is RF's
Tweedie identity. Our reimplementation lives in
\verb|src/samplers/flowdps_rf.py| in the project repository.

\paragraph{Baselines.}
\begin{itemize}[leftmargin=*,itemsep=2pt]
  \item \textbf{Wiener filter:} per-channel 2-D Wiener deconvolution in
    the Fourier domain, with noise-to-signal scalar $K = 10^{-3}$ for
    $\sigma_n = 0$ and $K = 5 \cdot 10^{-2}$ for $\sigma_n = 0.05$ (tuned
    on a 2-image validation).
  \item \textbf{DnCNN$+$PnP-ADMM:} PnP-ADMM in the Fourier domain with
    deepinv's pretrained 3-channel blind Gaussian DnCNN (668\,k
    parameters) as the denoiser. $24$ outer iterations, $\rho = 2.0$,
    descending denoiser-$\sigma$ schedule with $\sigma_\text{max} =
    \max(0.10, 4\sigma_n + 0.05)$ and $\sigma_\text{min} =
    \max(0.02, \sigma_n)$.
\end{itemize}

\paragraph{Implementation details.}
All experiments run on a single \textbf{NVIDIA RTX 3060 12\,GB} GPU.
Models load via Hugging Face's \texttt{diffusers} 0.30.1 on PyTorch
2.4.1$+$cu121. The CelebA-HQ-256 DDPM checkpoint is
\texttt{google/ddpm-ema-celebahq-256} (113\,M parameters); the
Rectified Flow checkpoint is the public CelebA-HQ-256 release from the
\texttt{gnobitab/RectifiedFlow} repository (65.5\,M parameters, NCSN$++$
architecture, JIT-compiled \texttt{upfirdn2d} CUDA op). LPIPS uses the
AlexNet backbone via the standard \texttt{lpips} package. Wall-clock
timings use CUDA events.

% =====================================================================
\section{Experimental Design}\label{sec:expdesign}

\paragraph{Dataset.}
We use 200 CelebA-HQ images from the validation split of the public
\texttt{korexyz/celeba-hq-256x256} dataset on Hugging Face. The first
50 images are used for the main experimental grid and the robustness
study. The 200-image set size is set by our compute envelope on the
RTX 3060 12\,GB; we disclose this explicitly so per-cell error bars
reflect this sample size.

\begin{table}[h]
  \centering
  \caption{Experimental configuration matrix.}
  \label{tab:config}
  \begin{tabular}{lll}
    \toprule
    Parameter & Values & Cells \\
    \midrule
    Gaussian blur $\sigma_b$ & $\{1.5, 3.0, 5.0\}$ & 3 \\
    Additive noise $\sigma_n$ & $\{0.0, 0.05\}$ & 2 \\
    NFE budget & $\{25, 50, 100\}$ & 3 \\
    Method & \{Pixel-DPS, FlowDPS-on-RF\} & 2 \\
    \midrule
    Total main-grid conditions & $3 \cdot 2 \cdot 3 \cdot 2 = 36$ & — \\
    Images per cell & $50$ & — \\
    Total reconstructions (main grid) & $1{,}800$ & — \\
    \midrule
    Robustness: motion length $L$ & $\{15, 25, 35\}$ & 3 \\
    Robustness: motion angle $\theta$ & $\{0^\circ, 45^\circ\}$ & 2 \\
    Robustness: NFE & $50$ (fixed) & — \\
    Total robustness reconstructions & $3 \cdot 2 \cdot 2 \cdot 2 \cdot 50 = 1{,}200$ & — \\
    \bottomrule
  \end{tabular}
\end{table}

\paragraph{Hyperparameter tuning.}
Guidance scale $\zeta$ tuned on a 2-image held-out subset (EXP-001 for
Pixel-DPS, EXP-006/007 for FlowDPS-on-RF). Locked values for the
\textbf{v1} baselines: $\zeta = 10$ for Pixel-DPS, $\zeta = 100$ for
FlowDPS-on-RF. Both are monotonic-best in their respective ranges;
values past these begin to hurt quality. The eight
publication-improvement variants in \S\ref{sec:ablations} each carry
their own re-tuned $\zeta$ (or $\zeta$-schedule) chosen on a disjoint
held-out subset (images 50--54), to avoid leakage into the
50-image evaluation set used for the significance tests.

\paragraph{Metrics.}
We report PSNR (dB, on the $[0, 1]$ range), SSIM (a PyTorch-native
implementation matching the conventional Gaussian-windowed formula),
and LPIPS (AlexNet backbone via the standard \texttt{lpips} package).
Higher is better for PSNR / SSIM; lower for LPIPS. Computational
efficiency is reported as wall-clock seconds per image (CUDA-event
timed, excluding I/O) and the number of neural function evaluations
(NFE) per reconstruction. We omit Kernel Inception Distance (KID),
which the proposal listed: distributional metrics over 50 samples per
cell are too noisy to be informative, and our paired-image PSNR / SSIM
/ LPIPS already capture the same quality axes.

\paragraph{Baselines.}
To contextualize the DPS / FlowDPS comparison we include two non-
diffusion baselines:
\begin{itemize}[leftmargin=*,itemsep=2pt]
  \item \textbf{Wiener filter}: per-channel 2-D Wiener deconvolution
    in the Fourier domain, with noise-to-signal scalar $K = 10^{-3}$
    for $\sigma_n = 0$ and $K = 5 \cdot 10^{-2}$ for $\sigma_n = 0.05$.
    Evaluated on the 6-cell Gaussian grid (300 reconstructions,
    $\sim$16\,ms per image on CPU).
  \item \textbf{DnCNN$+$PnP-ADMM}: Plug-and-Play ADMM in the Fourier
    domain with deepinv's pretrained 3-channel blind Gaussian DnCNN
    (668\,k parameters)~\cite{zhang2017dncnn,chan2017pnp}. 24 outer
    iterations, $\rho = 2.0$, descending denoiser-$\sigma$ schedule
    with $\sigma_\text{max} = \max(0.10, 4\sigma_n + 0.05)$ and
    $\sigma_\text{min} = \max(0.02, \sigma_n)$. 300 reconstructions,
    $\sim$0.3\,s per image on the same RTX 3060.
\end{itemize}

The proposal also listed DDRM~\cite{kawar2022ddrm} as a planned
baseline. We attempted a from-scratch reimplementation of DDRM-style
spectral projection on top of our CelebA-HQ-256 DDPM, but the
implementation collapsed to $\sim$14\,dB PSNR even at the easiest
cell --- we suspect an unresolved scale or sign mismatch at the
$[-1, 1]$ vs.\ $[0, 1]$ FFT boundary. To preserve scientific honesty
we report this attempt as a parked experiment (the code is preserved
in the project repository for future investigation) and use
DnCNN$+$PnP-ADMM as the supervised-classical reference in lieu of
DDRM. We also drop the unconditional-DDPM ``lower bound'' baseline
the proposal mentioned --- it would simply measure how far a random
sample from the prior is from the ground truth, which the other
baselines already bracket from above.

% =====================================================================
\section{Results and Findings}\label{sec:results}

\subsection{Main grid — matched Gaussian blur}


Table~\ref{tab:main_nfe100} shows per-cell mean PSNR/SSIM/LPIPS at
NFE${=}100$ for both methods; the full table including NFE${=}25,50$
is in the supplementary material.

\begin{table}[h]
  \centering
  \caption{Main-grid PSNR / SSIM / LPIPS at NFE${=}100$ (mean $\pm$ std over 50
    images per cell; best per cell in bold).}
  \label{tab:main_nfe100}
  \begin{tabular}{llrrlrr}
    \toprule
    Method & $\sigma_b$ & $\sigma_n$ & PSNR (dB) & SSIM & LPIPS & s/img \\
    \midrule
    Pixel-DPS     & 1.5 & 0.00 & \textbf{28.61 $\pm$ 3.18} & \textbf{0.821} & \textbf{0.084} & 19.0 \\
    Pixel-DPS     & 1.5 & 0.05 & \textbf{26.21 $\pm$ 2.39} & \textbf{0.686} & \textbf{0.136} & 17.9 \\
    Pixel-DPS     & 3.0 & 0.00 & \textbf{27.15 $\pm$ 2.97} & \textbf{0.775} & \textbf{0.117} & 18.0 \\
    Pixel-DPS     & 3.0 & 0.05 & \textbf{26.25 $\pm$ 2.64} & \textbf{0.741} & \textbf{0.146} & 17.9 \\
    Pixel-DPS     & 5.0 & 0.00 & \textbf{26.21 $\pm$ 2.72} & \textbf{0.746} & \textbf{0.155} & 17.9 \\
    Pixel-DPS     & 5.0 & 0.05 & \textbf{25.33 $\pm$ 2.46} & \textbf{0.721} & \textbf{0.190} & 18.1 \\
    \midrule
    FlowDPS-on-RF & 1.5 & 0.00 & 26.23 $\pm$ 2.66 & 0.737 & 0.153 & 23.5 \\
    FlowDPS-on-RF & 1.5 & 0.05 & 23.65 $\pm$ 1.86 & 0.570 & 0.268 & 23.5 \\
    FlowDPS-on-RF & 3.0 & 0.00 & 24.26 $\pm$ 2.41 & 0.667 & 0.170 & 23.7 \\
    FlowDPS-on-RF & 3.0 & 0.05 & 22.58 $\pm$ 1.96 & 0.576 & 0.234 & 23.9 \\
    FlowDPS-on-RF & 5.0 & 0.00 & 22.48 $\pm$ 2.23 & 0.601 & 0.200 & 23.8 \\
    FlowDPS-on-RF & 5.0 & 0.05 & 21.13 $\pm$ 1.87 & 0.537 & 0.251 & 24.0 \\
    \bottomrule
  \end{tabular}
\end{table}

\begin{figure}[h]
  \centering
  \includegraphics[width=\textwidth]{psnr_vs_nfe}
  \caption{PSNR vs NFE across all six (matched) degradation conditions.
    FlowDPS-on-RF saturates immediately; Pixel-DPS keeps improving.}
  \label{fig:psnr_vs_nfe}
\end{figure}

\begin{figure}[h]
  \centering
  \includegraphics[width=0.7\textwidth]{psnr_delta_heatmap}
  \caption{Per-cell PSNR advantage of Pixel-DPS over FlowDPS-on-RF.
    The gap widens with task difficulty, from $+1.4$\,dB at the easiest
    cell to $+4.2$\,dB at the hardest.}
  \label{fig:psnr_delta_heatmap}
\end{figure}

\begin{figure}[h]
  \centering
  \begin{subfigure}[t]{0.48\textwidth}
    \centering
    \includegraphics[width=\textwidth]{efficiency_frontier}
    \caption{Quality--efficiency frontier across all 18 cells
      (NFE-annotated). Pixel-DPS dominates at every NFE budget.}
    \label{fig:efficiency_frontier}
  \end{subfigure}
  \hfill
  \begin{subfigure}[t]{0.48\textwidth}
    \centering
    \includegraphics[width=\textwidth]{lpips_vs_psnr}
    \caption{LPIPS vs PSNR scatter. Light points $=$ per-image; dark
      outlined $=$ cell means. The two clouds barely overlap.}
    \label{fig:lpips_vs_psnr}
  \end{subfigure}
  \caption{Two cross-cell views of the matched-grid comparison.}
\end{figure}

\paragraph{Headline findings (matched Gaussian).}
\begin{enumerate}[leftmargin=*,itemsep=2pt]
  \item \textbf{Pixel-DPS dominates at every cell}, by $+1.4$ to
    $+4.2$\,dB PSNR. The gap grows with $\sigma_b$ and with
    $\sigma_n = 0.05$ vs.~$0.0$.
  \item \textbf{FlowDPS-on-RF saturates with NFE}: at
    $(\sigma_b{=}1.5, \sigma_n{=}0)$ the PSNR averages
    $26.08, 26.08, 26.23$\,dB across NFE${\in}\{25, 50, 100\}$.
    Pixel-DPS at the same cell: $27.51, 28.34, 28.61$\,dB — monotonic.
  \item \textbf{Pixel-DPS is also faster per image} on this hardware
    ($\sim$18\,s vs $\sim$24\,s at NFE${=}100$), reversing the
    proposal's hypothesized flow-speed advantage.
\end{enumerate}

\subsection{Operator-mismatch robustness study}

The true forward operator is $\text{motion}(L, \theta)$ but the sampler's
likelihood assumes Gaussian with fixed $\sigma_b{=}3.0$. 50 images per
cell, $3 \times 2 \times 2 = 12$ cells per method, NFE${=}50$.

\begin{figure}[h]
  \centering
  \includegraphics[width=0.9\textwidth]{matched_vs_mismatched}
  \caption{Operator-mismatch cost ($\sigma_n{=}0.05$, NFE${=}50$).
    Pixel-DPS drops $\sim$5\,dB from matched Gaussian to severe motion
    ($L{=}35$); FlowDPS-on-RF drops only $\sim$2.5\,dB.}
  \label{fig:matched_vs_mismatched}
\end{figure}

\begin{figure}[h]
  \centering
  \includegraphics[width=\textwidth]{robustness_psnr_heatmap}
  \caption{Per-cell mean PSNR under operator mismatch (true forward $=$
    motion blur, sampler assumes Gaussian with fixed $\sigma{=}3.0$).
    Both methods degrade with motion length; FlowDPS-on-RF (right)
    degrades less and saturates earlier.}
  \label{fig:robustness_heatmap}
\end{figure}

\paragraph{Headline finding (robustness).}
\textbf{FlowDPS-on-RF is more robust to operator mismatch than
Pixel-DPS.} Going from matched Gaussian ($\sigma_b{=}3, \sigma_n{=}0.05,$
NFE${=}50$) to the hardest motion case ($L{=}35, \theta{=}45^\circ$):
\begin{itemize}[leftmargin=*,itemsep=2pt]
  \item Pixel-DPS: $25.72 \to 20.86$\,dB ($\boldsymbol{-4.86}$\,dB).
  \item FlowDPS-on-RF: $22.40 \to 19.92$\,dB ($\boldsymbol{-2.48}$\,dB).
\end{itemize}
The matched-condition $+3.3$\,dB Pixel-DPS advantage shrinks to
$+0.9$\,dB at the hardest mismatched cell. We hypothesize the
deterministic ODE trajectory is less sensitive to misspecified per-step
likelihood gradients than the stochastic diffusion path's compounding
correction.

\subsection{Classical and supervised baselines}

\begin{table}[h]
  \centering
  \caption{Four-method PSNR comparison at $\sigma_n{=}0.05$, NFE${=}100$
    where applicable. Per-method mean $\pm$ std over 50 images per cell;
    best per cell in bold.}
  \label{tab:baselines}
  \begin{tabular}{llll>{\bfseries}l}
    \toprule
    $\sigma_b$ & Wiener & DnCNN$+$PnP-ADMM & FlowDPS-on-RF & Pixel-DPS \\
    \midrule
    1.5 & 24.28 $\pm$ 1.29 & 24.61 $\pm$ 1.21 & 23.65 $\pm$ 1.86 & 26.21 $\pm$ 2.39 \\
    3.0 & 24.03 $\pm$ 1.73 & 24.31 $\pm$ 1.57 & 22.58 $\pm$ 1.96 & 26.25 $\pm$ 2.64 \\
    5.0 & 22.50 $\pm$ 1.70 & 22.81 $\pm$ 1.54 & 21.13 $\pm$ 1.87 & 25.33 $\pm$ 2.46 \\
    \midrule
    s/img & 0.02 & 0.30 & 23.5 & 17.9 \\
    \bottomrule
  \end{tabular}
\end{table}

Three orders of magnitude in compute for roughly $2$\,dB in PSNR. The
DnCNN$+$PnP baseline beats our FlowDPS reimplementation in $5/6$
$\sigma_n{=}0.05$ cells while running $\sim$60$\times$ faster than DPS.
Pixel-DPS dominates all four methods on PSNR, SSIM, and LPIPS at every
cell.

\begin{figure}[h]
  \centering
  \includegraphics[width=\textwidth]{qualitative_grid}
  \caption{Qualitative comparison on six representative CelebA-HQ-256 test
    images at $\sigma_b{=}3.0$, $\sigma_n{=}0.05$, NFE${=}50$. Columns:
    clean ground truth, degraded measurement, Wiener filter,
    DnCNN$+$PnP-ADMM, FlowDPS-on-RF, and Pixel-DPS. Numbers below each
    reconstruction are per-image PSNR (dB). Wiener and DnCNN$+$PnP are
    smooth but oversmooth identity-relevant detail; FlowDPS-on-RF
    sometimes drifts to a plausible-but-wrong face; Pixel-DPS most
    consistently recovers identity-preserving facial structure.}
  \label{fig:qualitative_grid}
\end{figure}

\subsection{Posterior diversity vs.\ point accuracy}\label{sec:diversity}

A high PSNR on a single sample is not the same as a good posterior
sampler. The proposal brief (and the standard inverse-problem
literature) emphasizes the distinction between (i) the \emph{best}
sample from a posterior, (ii) the posterior \emph{mean}, and (iii) the
sample \emph{diversity}. To probe this for Pixel-DPS and FlowDPS-on-RF
we draw $N{=}8$ samples per image with different sampler seeds (the
measurement $y$ and its noise are held fixed across seeds) on three
representative images at $\sigma_b{=}3.0$, $\sigma_n{=}0.05$, NFE${=}50$.

We report four quantities per image $\times$ method:
\begin{itemize}[leftmargin=*,itemsep=2pt]
  \item Single-sample PSNR (seed 100, i.e.\ what we have been calling
    ``the reconstruction'' throughout this paper).
  \item Mean-of-$N$ PSNR --- the PSNR of the per-pixel mean of the 8
    sampled reconstructions. This estimates the posterior mean.
  \item Best-of-$N$ PSNR --- the maximum per-image PSNR across the 8
    samples (an oracle that knows the ground truth).
  \item Pairwise mean LPIPS over $\binom{N}{2}{=}28$ sample pairs ---
    a perceptual measure of posterior sample diversity. Higher means
    the sampler explores more of the posterior; lower means it
    collapses to a near-deterministic mode.
\end{itemize}

\begin{table}[h]
  \centering
  \caption{Posterior-diversity study at $\sigma_b{=}3.0$, $\sigma_n{=}0.05$,
    NFE${=}50$, $N{=}8$ samples per (image, method). Mean-of-$N$ is the
    PSNR of the per-pixel sample mean. Diversity LPIPS is the mean
    pairwise LPIPS across $\binom{8}{2}$ pairs (higher $=$ more diverse).}
  \label{tab:diversity}
  \begin{tabular}{llrrrrr}
    \toprule
    Image & Method & Single & Mean-of-$N$ & Best-of-$N$ & Div.\ LPIPS & std map (mean) \\
    \midrule
    0  & Pixel-DPS     & 31.25 & 32.33 & 31.58 & 0.043 & 0.009 \\
    0  & FlowDPS-on-RF & 27.46 & 29.38 & 27.72 & \textbf{0.233} & \textbf{0.032} \\
    7  & Pixel-DPS     & 28.23 & 29.14 & 28.37 & 0.057 & 0.012 \\
    7  & FlowDPS-on-RF & 24.85 & 27.09 & 25.20 & \textbf{0.151} & \textbf{0.031} \\
    25 & Pixel-DPS     & 27.73 & 28.42 & 27.73 & 0.050 & 0.014 \\
    25 & FlowDPS-on-RF & 24.41 & 26.68 & 25.21 & \textbf{0.142} & \textbf{0.030} \\
    \midrule
    \textit{average}  & Pixel-DPS     & 29.07 & 29.96 & 29.23 & 0.050 & 0.012 \\
    \textit{average}  & FlowDPS-on-RF & 25.57 & 27.72 & 26.04 & \textbf{0.175} & \textbf{0.031} \\
    \bottomrule
  \end{tabular}
\end{table}

\begin{figure}[h]
  \centering
  \includegraphics[width=\textwidth]{posterior_diversity}
  \caption{Posterior diversity. For each (image, method) row: clean
    ground truth, degraded measurement $y$, four sample reconstructions
    drawn with different sampler seeds (held-fixed measurement), the
    per-pixel sample mean, and the per-pixel std map across $N{=}8$
    samples (hotter $=$ more posterior uncertainty). Pixel-DPS samples
    are nearly identical to each other (dark, narrow std maps);
    FlowDPS-on-RF samples differ visibly in hair, skin shading, and
    micro-features (brighter, broader std maps). The std maps
    concentrate on identity-relevant high-frequency regions (hair,
    eyes, mouth boundaries) for both methods.}
  \label{fig:diversity}
\end{figure}

\paragraph{Headline findings (diversity).}
\begin{enumerate}[leftmargin=*,itemsep=2pt]
  \item \textbf{Pixel-DPS is near-deterministic across seeds.} Pairwise
    LPIPS across 8 samples is $0.05$ on average, and the per-pixel std
    map averages $0.012$ on the $[0, 1]$ scale --- the sampler is
    essentially collapsing to a single posterior mode. This is the
    failure mode the inverse-problem literature warns about: high
    single-sample PSNR can coexist with poor posterior coverage.
  \item \textbf{FlowDPS-on-RF samples are 3--4$\times$ more diverse}
    (LPIPS $0.175$, std mean $0.031$). Different seeds yield visibly
    different hair, skin shading, and micro-features (Figure~\ref{fig:diversity}).
  \item \textbf{The posterior mean strictly improves PSNR} for both
    methods, but the gain is much larger for FlowDPS-on-RF (mean-of-8
    minus single $= {+}2.15$\,dB) than for Pixel-DPS ($+0.89$\,dB).
    This is the classic Bayesian gain from averaging out independent
    sample noise; FlowDPS-on-RF benefits more precisely because its
    individual samples are noisier (more diverse).
  \item \textbf{Best-of-$N$ is barely above single-sample for either
    method.} The 8 samples cluster too tightly around the modal
    estimate to yield large best-case improvements. A best-of-1000
    study or a learned sample selector might change this.
\end{enumerate}

Reading these together: the matched-Gaussian PSNR / SSIM / LPIPS
advantage Pixel-DPS holds in Tables~\ref{tab:main_nfe100}
and~\ref{tab:baselines} is partly a single-sample-point estimate
advantage rather than a true posterior-sampling advantage. From a
Bayesian inverse-problem perspective, FlowDPS-on-RF is
\emph{producing} a posterior; Pixel-DPS is producing a
\emph{maximum-a-posteriori-like point estimate}. Which is preferable
depends on whether downstream use cares about modal accuracy or
uncertainty awareness.

\subsection{Statistical significance}\label{sec:significance}

To assess whether the per-cell PSNR differences in
Tables~\ref{tab:main_nfe100}/\ref{tab:baselines} are statistically
distinguishable from noise across the 50-image samples we used, we
run paired t-tests and Wilcoxon signed-rank tests on the per-image
PSNR vectors for each (Pixel-DPS, other-method, $\sigma_b$, $\sigma_n$)
combination at NFE${=}100$. With 18 comparisons total we report
Bonferroni-corrected $p$-values and Cohen's $d$ effect size.

\begin{table}[h]
  \centering
  \caption{Paired-t test of Pixel-DPS vs each other method (NFE${=}100$,
    matched $\sigma_b$, $\sigma_n$, $n{=}50$ images). Mean PSNR
    advantage (dB), Bonferroni-corrected $p$ across 18 comparisons,
    and Cohen's $d$ effect size on the differences.}
  \label{tab:significance}
  \begin{tabular}{llrrr}
    \toprule
    Comparison & ($\sigma_b$, $\sigma_n$) & $\Delta$PSNR (dB) & $p$ (Bonferroni) & Cohen's $d$ \\
    \midrule
    vs FlowDPS-on-RF & (1.5, 0.00) & $+2.39$ & $1.8 \times 10^{-17}$ & 1.98 \\
    vs FlowDPS-on-RF & (1.5, 0.05) & $+2.56$ & $1.4 \times 10^{-21}$ & 2.49 \\
    vs FlowDPS-on-RF & (3.0, 0.00) & $+2.89$ & $1.2 \times 10^{-19}$ & 2.24 \\
    vs FlowDPS-on-RF & (3.0, 0.05) & $+3.67$ & $3.1 \times 10^{-25}$ & 3.02 \\
    vs FlowDPS-on-RF & (5.0, 0.00) & $+3.73$ & $6.8 \times 10^{-24}$ & 2.82 \\
    vs FlowDPS-on-RF & (5.0, 0.05) & $+4.20$ & $3.2 \times 10^{-28}$ & 3.52 \\
    \midrule
    vs Wiener        & (1.5, 0.00) & $+6.39$ & $8.1 \times 10^{-20}$ & 2.26 \\
    vs Wiener        & (1.5, 0.05) & $+1.93$ & $2.5 \times 10^{-13}$ & 1.53 \\
    vs Wiener        & (3.0, 0.00) & $+7.21$ & $3.4 \times 10^{-23}$ & 2.72 \\
    vs Wiener        & (3.0, 0.05) & $+2.21$ & $5.1 \times 10^{-17}$ & 1.93 \\
    vs Wiener        & (5.0, 0.00) & $+8.25$ & $2.2 \times 10^{-26}$ & 3.21 \\
    vs Wiener        & (5.0, 0.05) & $+2.83$ & $3.9 \times 10^{-21}$ & 2.43 \\
    \midrule
    vs DnCNN$+$PnP   & (1.5, 0.00) & $+0.16$ & $\mathbf{1.0}$ (\textbf{n.s.}) & 0.21 \\
    vs DnCNN$+$PnP   & (1.5, 0.05) & $+1.60$ & $4.9 \times 10^{-10}$ & 1.21 \\
    vs DnCNN$+$PnP   & (3.0, 0.00) & $+1.50$ & $2.2 \times 10^{-12}$ & 1.44 \\
    vs DnCNN$+$PnP   & (3.0, 0.05) & $+1.94$ & $3.7 \times 10^{-14}$ & 1.62 \\
    vs DnCNN$+$PnP   & (5.0, 0.00) & $+2.77$ & $1.3 \times 10^{-19}$ & 2.24 \\
    vs DnCNN$+$PnP   & (5.0, 0.05) & $+2.52$ & $2.5 \times 10^{-19}$ & 2.20 \\
    \bottomrule
  \end{tabular}
\end{table}

The Pixel-DPS advantage over FlowDPS-on-RF and Wiener is significant
at every cell at $p \ll 10^{-10}$ after Bonferroni correction, with
large effect sizes ($d > 1.5$ everywhere, $> 2.5$ in 8 of 12 cells).

The comparison against DnCNN$+$PnP-ADMM is more nuanced: at the
\emph{easiest} cell ($\sigma_b{=}1.5$, $\sigma_n{=}0$), Pixel-DPS's
$+0.16$\,dB margin is \emph{not statistically significant}
($p{=}0.148$ raw, $p{=}1.0$ Bonferroni). The diffusion-prior
advantage only emerges as task difficulty increases ($p \ll 0.001$
for the five harder cells, with the advantage growing to $+2.8$\,dB
at the hardest). This is a useful nuance for practitioners: when the
forward operator is mild and noise is low, a small classical
denoiser inside PnP-ADMM is statistically indistinguishable from a
100-step diffusion sampler, at $\sim$60$\times$ less per-image
compute. Diffusion priors earn their compute budget on the
harder cells.

% =====================================================================
\subsection{Publication-improvement ablations}\label{sec:ablations}

Beyond the v1 baselines reported above, we explored \emph{eight}
candidate improvements to either sampler grouped into three tiers
(A/B/C) plus a follow-up round (R2). Each variant defines a new method
name (e.g.\ \texttt{flowdps\_rf\_v2}) so the original v1 rows in our
results CSV remain untouched as the ``before'' reference. The promotion
criterion was strict: a variant passes the gate only if (i) its mean
PSNR difference over the v1 baseline is \emph{positive} AND
(ii) Bonferroni-corrected $p < 0.05$ in $\geq 4/6$
$(\sigma_b, \sigma_n)$ cells, at NFE${=}100$, $n{=}50$ images per cell.
\textbf{Of the eight candidates only one passed the gate}; the rest
are documented here as instructive negative results rather than
discarded. The complete scorecard appears in
Table~\ref{tab:ablations}.

\begin{table}[h]
  \centering
  \caption{Round-1 (Tier A/B/C) and Round-2 publication-improvement
    candidates against their v1 baseline. ``Wins/6'' counts cells with
    positive $\Delta$ AND Bonferroni $p<0.05$. Only
    \texttt{flowdps\_rf\_v2} passes the gate.}
  \label{tab:ablations}
  \begin{tabular}{llrrl}
    \toprule
    Method & Tier & Mean $\Delta$ (dB) & Wins / 6 & Verdict \\
    \midrule
    \texttt{flowdps\_rf\_v2}   & A & $\boldsymbol{+0.77}$ & $\boldsymbol{6/6}$ & \textbf{PASSES} \\
    \texttt{pixel\_dps\_v2}    & A & $-0.66$ & $0/6$ & fail \\
    \texttt{pixel\_dps\_spectral}\textsuperscript{*}  & B & $-0.18$ & $0/12$ & fail \\
    \texttt{flowdps\_rf\_spectral}\textsuperscript{*} & B & $-0.20$ & $0/12$ & fail \\
    \texttt{pixel\_dps\_pigdm}  & C & $-13.63$ & $0/6$ & fail \\
    \texttt{flowdps\_rf\_pigdm} & C & $-9.81$ & $0/6$ & fail \\
    \texttt{pixel\_dps\_sched}  & R2 & $-0.46$ & $2/6$ & fail \\
    \texttt{flowdps\_rf\_heun}  & R2 & $+0.53$ vs v1 \,/\, $-0.25$ vs v2 & $0/6$ vs v2 & fail (dominated by v2) \\
    \bottomrule
  \end{tabular}
  \\[0.4em]
  \footnotesize \textsuperscript{*}Spectral methods are evaluated on
  the 12-cell motion-blur robustness study (their design target),
  not the matched-Gaussian main grid.
\end{table}

\paragraph{Tier-A: EMA weights $+$ $\zeta$ schedule (1 win, 1 fail).}
The Liu~2023 RF checkpoint ships 644 EMA \texttt{shadow\_params} for a
645-parameter network. The missing parameter is the Gaussian
Fourier-features projection \texttt{all\_modules.0.W}, registered with
\texttt{requires\_grad=False} and therefore skipped by
\texttt{gnobitab}'s EMA tracker. \texttt{flowdps\_rf\_v2} aligns the
644 trainable params to the shadow\_params \emph{and} separately
copies the frozen \texttt{all\_modules.0.W} from the raw \texttt{model}
state-dict (omitting the second step leaves the Fourier projection at
its random initialization and yields catastrophic 7.76\,dB PSNR). It
then replaces the scalar $\zeta{=}100$ with a power-law ramp
$\zeta(i) = 200 \cdot (i/N)^{0.5}$, which puts more guidance toward
the data side of the trajectory where the FlowDPS-RF v1 sampler
saturates earliest (Figure~\ref{fig:psnr_vs_nfe}). Mean improvement
over v1: $\boldsymbol{+0.77}$\,dB PSNR, with Cohen's $d$ ranging
$1.39$--$3.17$ across the six cells (a very large effect by
conventional thresholds). The companion attempt \texttt{pixel\_dps\_v2}
(scalar $\zeta = 30$, the analogous Pixel-DPS tweak) failed in the
opposite direction: it over-corrects on easy cells
($-0.93$\,dB at $\sigma_b{=}1.5,\sigma_n{=}0$) and gives marginal
gains only on noisier cells, netting $-0.66$\,dB on average. A single
global scalar $\zeta$ cannot satisfy both easy and hard cells in
Pixel-DPS; this directly motivated the Round-2 \texttt{pixel\_dps\_sched}
attempt below.

\paragraph{Tier-B: spectral-aware likelihood (2 fails).}
Under operator mismatch the residual energy in the bands the assumed
operator cannot probe is dominated by measurement noise, suggesting a
frequency-dependent weight
$W(f) = 1/(1 + \alpha \cdot \text{relu}(\epsilon - |H(f)|))$
that downweights residuals where the assumed (Gaussian) operator's
spectral magnitude is low. We tested both
\texttt{pixel\_dps\_spectral} and \texttt{flowdps\_rf\_spectral}
on the motion-blur robustness study (12 cells $\times$ 50 imgs $=$
600 reconstructions per method). Both fail (mean $-0.18$\,dB and
$-0.20$\,dB respectively; 10/12 and 5/12 cells Bonferroni-significant
but every significant cell is in the negative direction). The failure
is structural rather than tuning-induced: the true motion-blur
operator has support in different frequency bands than the assumed
Gaussian, so suppressing residuals in the latter's null space also
throws away informative signal that the true operator preserved.
The noise-floor heuristic was fighting the wrong axis.

\paragraph{Tier-C: Tweedie-corrected ($\Pi$-GDM-style) likelihood (2 fails, with a subtle bug).}
The proper posterior-aware likelihood under Tweedie's clean-image
estimate (covariance $r_t \cdot I$ where
$r_t = (1-\bar\alpha_t)/\bar\alpha_t$ for diffusion or $(1-t)^2$
for rectified flow) is Gaussian with covariance
$\sigma_n^2 I + r_t \cdot A A^\top$, diagonal in the FFT basis. Our
\emph{first} attempt used a mean-normalized inverse-variance weight
combined with the existing DPS L2-\emph{norm} loss; this collapsed
PSNR to $\sim$10\,dB across all 18 cells. The diagnosis: $\Pi$-GDM is
by derivation an L2-\emph{squared} whitened-residual likelihood, and
mixing the two conventions scrambles the per-frequency gradient
scaling. After correcting to un-normalized $W$ $+$ L2-norm $+$ per-cell
OTF (un-normalized $W$ preserves the correct relative scaling across
frequencies; L2-norm's $1/\text{loss}$ factor bounds gradient
magnitude across the ${\sim}4 \times 10^5{\times}$ swing of $r_t$
over the schedule), and tuning $\zeta{=}100$ via a held-out 2-image
sweep, the corrected implementation still fails the gate
(\texttt{pixel\_dps\_pigdm}: $-13.6$\,dB; \texttt{flowdps\_rf\_pigdm}:
$-9.8$\,dB). The 18-cell grid splits sharply by noise level: the
$\sigma_n{=}0$ cells collapse to 6--9\,dB because $W$ scales like
$1/\sigma_n$ at small $r_t$ and the $10^{-3}$ noise-floor clamp makes
the effective $\zeta$ explode, while the $\sigma_n{=}0.05$ cells
merely lose by $0.5$--$8.7$\,dB. A practical implementation would
require $\sigma_n$-dependent $\zeta$ scaling, which we leave to
follow-up work --- and which still does not appear to flip the sign
even when correctly tuned, suggesting $\Pi$-GDM-style corrections
benefit inpainting and super-resolution more than Gaussian
deblurring on this prior.

\paragraph{Round-2: $\zeta$-schedule (Pixel-DPS) and Heun integrator (FlowDPS-RF) (2 fails).}
Two follow-up attempts inspired by the Tier-A win.
\texttt{pixel\_dps\_sched} replaces the scalar $\zeta{=}10$ with the
best-of-eight schedule chosen on \emph{held-out} images 50--54
(\texttt{ramp(80, 0.5)}), the genuine Pixel-DPS analog of what made
\texttt{flowdps\_rf\_v2} win. On the full 50-image grid it fails the
gate (mean $-0.46$\,dB), with a clean cell-level split: it
\emph{wins} ${+}0.8$\,dB on both $\sigma_n{=}0.05$ cells but
\emph{loses} on every $\sigma_n{=}0$ cell. The optimal $\zeta$
scaling for Pixel-DPS apparently depends on measurement noise level,
not on timestep --- the same diagnosis as Tier-C in a different guise.
\texttt{flowdps\_rf\_heun} swaps the Euler integrator for a 2nd-order
Heun predictor-corrector on top of the Tier-A winning v2 config, with
the NFE budget preserved by halving the step count. The result is
more honest if subtle: \texttt{flowdps\_rf\_heun} technically passes
the significance gate against \emph{v1} ($+0.53$\,dB) because it
inherits most of v2's win, but against its actual base config v2 it
is a consistent $-0.25$\,dB regression across all six cells
($p < 10^{-19}$ in each cell). At a fixed NFE budget Heun spends two
velocity calls per step, so it runs only $N/2$ steps and thus invokes
the measurement-consistency update only $N/2$ times --- half as often
as Euler. The integration-accuracy gain does not pay for the loss of
guidance correction frequency on this problem.

\paragraph{What the failed ablations tell us.}
Read together, the seven failed attempts share a common pattern: each
extra mechanism (higher fixed $\zeta$, spectral-domain reweighting,
$\Pi$-GDM correction, time-varying $\zeta$, higher-order integrator)
imposes a single global hyperparameter on a problem where the optimal
setting actually varies along an axis the mechanism does not model
(cell difficulty, true vs.\ assumed operator support, $\sigma_n$, step
count). The lone Tier-A win, \texttt{flowdps\_rf\_v2}, succeeds
precisely because it introduces a \emph{time-varying} schedule (more
guidance toward the data side) on a sampler whose existing failure
mode --- saturating early in NFE (\S\ref{sec:results},
Figure~\ref{fig:psnr_vs_nfe}) --- benefits from increasing late-step
measurement consistency. None of the candidates beats Pixel-DPS in
absolute terms; Tier-A narrows the matched-Gaussian PSNR gap from
$+3.24$\,dB to roughly $+1.4$\,dB (e.g.\ 28.61 vs.\ 27.35 at
$\sigma_b{=}1.5, \sigma_n{=}0$, NFE${=}100$) but the directional
finding from \S\ref{sec:results} holds.

% =====================================================================
\section{Discussion}\label{sec:discussion}

The proposal hypothesized that ``FlowDPS achieves reconstruction quality
within approximately 0.5--1.1\,dB PSNR of DPS while offering a substantial
$2$--$5\times$ reduction in computational cost.'' Our measurements
\emph{contradict both halves} of this claim on the CelebA-HQ-256
pixel-space comparison: (i) the PSNR gap is $+1.4$ to $+4.2$\,dB in
Pixel-DPS's favor, $2$--$4\times$ wider than predicted; and (ii)
Pixel-DPS is also faster per image at NFE${=}100$ on the RTX 3060,
reversing the speed ordering. The flow model's expected efficiency
advantage does not manifest in this regime; the RF velocity-field
NCSN$++$ forward pass is more expensive per step than the DDPM's
$\epsilon$ prediction, and at NFE${=}100$ both methods are far from any
flow-saturation regime where the straight-line trajectory would pay off
relative to a stochastic diffusion path.

However, the comparison \emph{reverses} under operator mismatch.
Pixel-DPS degrades by $\sim$5\,dB PSNR going from matched Gaussian to
severe motion-blur measurements, versus only $\sim$2.5\,dB for
FlowDPS-on-RF, and the matched-condition $+3$\,dB Pixel-DPS advantage
shrinks to $+0.9$\,dB at the hardest mismatched cell. This is the
\emph{opposite} of the proposal's hypothesis (which we had borrowed from
prior intuition that stochastic diffusion offers ``implicit
regularization'' against operator misspecification).

\paragraph{Point estimate vs.\ posterior sampling.}
The posterior-diversity study in \S\ref{sec:diversity} adds a third
axis to the comparison that the headline PSNR / SSIM / LPIPS numbers
hide. Pixel-DPS's pairwise sample LPIPS across 8 seeds is $\sim$0.05
on average --- the sampler is near-deterministic, collapsing to a
single modal estimate. FlowDPS-on-RF's pairwise sample LPIPS is
$\sim$0.18, a 3.5$\times$ broader posterior. Read together with the
matched-Gaussian PSNR table, this means \emph{Pixel-DPS's single-sample
PSNR advantage is partly a point-estimate advantage}, not a true
posterior-sampling advantage. The classical Bayesian-inverse-problems
critique --- that a sharp single sample with poor diversity is not the
same as a well-calibrated posterior --- applies here. From a downstream
perspective: if the application wants a single best-guess
reconstruction, Pixel-DPS dominates; if it wants samples that reflect
genuine measurement uncertainty (e.g.\ for downstream Bayesian
inference, uncertainty-aware decision making, or scientific imaging
where physically valid \emph{distributions} matter as much as image
quality), FlowDPS-on-RF is the better-calibrated sampler, despite its
worse point-estimate PSNR.

\paragraph{Gradient-based vs.\ gradient-free guidance.}
Both methods compared here are \emph{gradient-based}: they backpropagate
the measurement-likelihood gradient through the model at every step,
which roughly doubles the per-step compute relative to a forward-only
pass. Recent work~\cite{dou2026particle} proposes \emph{gradient-free}
constrained-particle methods that solve diffusion inverse problems
using only forward passes through the score / flow network. Comparing
gradient-based DPS against a gradient-free variant on the same
pixel-space DDPM backbone would isolate exactly how much of our
20-second-per-image cost is attributable to the autograd backward
pass, and is a clear direction for follow-up work.

\paragraph{Negative-result yield of the ablation study.}
The 1-pass-from-8 result of \S\ref{sec:ablations} is itself a finding.
Most of the candidate improvements --- spectral-domain reweighting,
$\Pi$-GDM-style covariance correction, higher-order integration ---
have textbook motivation but do not transfer cleanly to the
pixel-space CelebA-HQ-256 Gaussian-deblurring setting at the noise
levels and NFE budgets we test. The common failure mode is a global
hyperparameter masking a setting that should in fact vary along an
unmodeled axis (noise level, frequency band, or step count). Each
failed ablation is preserved as code in the project repository so
future work can revisit them with the axis-dependent tuning we
identified as the missing piece.

\paragraph{Limitations.}
(i)~50 images per cell — error bars are reported as $\pm\sigma$ in the
supplementary tables. (ii)~Single dataset (CelebA-HQ-256 faces).
(iii)~Single forward-operator family per study (Gaussian / motion blur).
(iv)~Our DDRM reimplementation did not reach competitive PSNR in the
available time (it collapsed to $\sim$14\,dB even at the easiest cell);
we report DnCNN$+$PnP-ADMM as the analogous classical-prior reference
instead. The parked DDRM code is preserved in the repository for future
investigation. (v)~All headline numbers for FlowDPS-on-RF in
\S\ref{sec:results} use the non-EMA \texttt{model} state-dict (matching
the v1 baseline); \S\ref{sec:ablations}'s \texttt{flowdps\_rf\_v2}
shows that correctly loading the EMA shadow\_params \emph{and} the
frozen Fourier-features projection adds $+0.77$\,dB to the
matched-Gaussian results, narrowing but not closing the
Pixel-DPS-vs-FlowDPS-RF gap.

\paragraph{Future work.}
A clean DDRM implementation would close one gap in our baseline coverage.
Investigating Kim~2025's per-step \texttt{step\_size} parameter or a
learned ODE-time schedule might close some of the Pixel-DPS vs
FlowDPS-on-RF matched-grid gap. Finally, repeating the comparison on
non-face datasets (LSUN-Church, DIV2K) would test how much of the
Pixel-DPS dominance is dataset-specific.

% =====================================================================
\section{Conclusion}\label{sec:conclusion}

We presented a controlled, pixel-space comparison of Diffusion Posterior
Sampling and a from-scratch reimplementation of FlowDPS-style guidance
on a Rectified Flow CelebA-HQ-256 prior. Pixel-DPS dominates the
matched-Gaussian comparison by $1.4$--$4.2$\,dB PSNR across an
$18$-condition grid, but FlowDPS-on-RF is more robust to operator
mismatch ($-2.5$\,dB drop vs.\ Pixel-DPS's $-5$\,dB drop under severe
motion-blur measurement mismatch). All four methods comfortably beat a
Wiener-filter baseline; DnCNN$+$PnP-ADMM sits in the efficient middle of
the Pareto frontier, beating FlowDPS-on-RF in $5/6$ noisy cells while
running $\sim$60$\times$ faster than DPS. We further evaluated eight
publication-improvement candidates against a strict Bonferroni-corrected
gate: only \texttt{flowdps\_rf\_v2} (correctly EMA-loaded weights $+$
power-law $\zeta$ ramp) passed, adding $+0.77$\,dB to FlowDPS-on-RF;
the seven failed candidates are documented as diagnosed negative
results, including a DPS-vs-$\Pi$-GDM L2-norm/squared convention
mismatch and a guidance-frequency-vs-integration-accuracy trade-off in
Heun. The findings give concrete guidance for choosing between these
methods under known versus uncertain forward operators: prefer DPS when
the operator is calibrated; consider flow-based posterior sampling when
the operator may be mis-specified.

% =====================================================================
% natbib author-year mode (forced by neurips_2024.sty) — each \bibitem
% must include the [Author(Year)] label so \citet / \citep can resolve.
\begin{thebibliography}{99}

\bibitem[Chung et~al.(2023)]{chung2023dps}
H.~Chung, J.~Kim, M.~T.~Mccann, M.~L.~Klasky, and J.~C.~Ye.
\emph{Diffusion Posterior Sampling for General Noisy Inverse Problems}.
ICLR 2023.

\bibitem[Kim et~al.(2025)]{kim2025flowdps}
J.~Kim, B.~S.~Kim, and J.~C.~Ye.
\emph{FlowDPS: Flow-Driven Posterior Sampling for Inverse Problems}.
ICCV 2025. arXiv:2503.08136.

\bibitem[Liu et~al.(2023)]{liu2023rectified}
X.~Liu, C.~Gong, and Q.~Liu.
\emph{Flow Straight and Fast: Learning to Generate and Transfer Data
with Rectified Flow}. ICLR 2023.

\bibitem[Lipman et~al.(2023)]{lipman2023flow}
Y.~Lipman, R.~T.~Q.~Chen, H.~Ben-Hamu, M.~Nickel, and M.~Le.
\emph{Flow Matching for Generative Modeling}. ICLR 2023.

\bibitem[Ho et~al.(2020)]{ho2020ddpm}
J.~Ho, A.~Jain, and P.~Abbeel.
\emph{Denoising Diffusion Probabilistic Models}. NeurIPS 2020.

\bibitem[Kawar et~al.(2022)]{kawar2022ddrm}
B.~Kawar, M.~Elad, S.~Ermon, and J.~Song.
\emph{Denoising Diffusion Restoration Models}. NeurIPS 2022.

\bibitem[Chan et~al.(2017)]{chan2017pnp}
S.~H.~Chan, X.~Wang, and O.~A.~Elgendy.
\emph{Plug-and-Play ADMM for Image Restoration: Fixed-Point Convergence
and Applications}. IEEE Trans.~Computational Imaging 2017.

\bibitem[Zhang et~al.(2017)]{zhang2017dncnn}
K.~Zhang, W.~Zuo, Y.~Chen, D.~Meng, and L.~Zhang.
\emph{Beyond a Gaussian Denoiser: Residual Learning of Deep CNN for
Image Denoising}. IEEE TIP 2017.

\bibitem[Esser et~al.(2024)]{esser2024sd3}
P.~Esser et~al.
\emph{Scaling Rectified Flow Transformers for High-Resolution Image
Synthesis}. arXiv:2403.03206, 2024.

\bibitem[Hendrycks and Dietterich(2019)]{hendrycks2019corrupt}
D.~Hendrycks and T.~Dietterich.
\emph{Benchmarking Neural Network Robustness to Common Corruptions and
Perturbations}. ICLR 2019.

\bibitem[Meanti et~al.(2025)]{meanti2025dim}
G.~Meanti et~al.
\emph{Unsupervised Imaging Inverse Problems with Diffusion Distribution
Matching}. ICCV 2025.

\bibitem[Dou et~al.(2026)]{dou2026particle}
H.~Dou et~al.
\emph{Constrained Particle Seeking: Solving Diffusion Inverse Problems
with Just Forward Passes}. arXiv:2603.01837, 2026.

\bibitem[Dasgupta et~al.(2026)]{dasgupta2026pcfm}
A.~Dasgupta et~al.
\emph{Solving Physics-Constrained Inverse Problems with Conditional
Flow Matching}. arXiv:2603.14135, 2026.

\end{thebibliography}

\end{document}
